\documentclass{article}
\usepackage{amsmath, amssymb, amsthm}
\usepackage{hyperref}
\usepackage{geometry}
\geometry{margin=1in}

\title{Erd\H{o}s Problem \#969}
\date{Last updated: 2025-08-31}
\author{Source: erdosproblems.com}

\begin{document}
\maketitle

\noindent\textbf{Status:} OPEN \quad \textbf{No prize}

\noindent\textbf{Tags:} number theory

\medskip
\noindent\textbf{Problem Statement:}

Let $Q(x)$ count the number of squarefree integers in $[1,x]$. Determine the order of magnitude in the error term in the asymptotic\[Q(x)=\frac{6}{\pi^2}x+E(x).\]

\bigskip
\noindent\textbf{Remarks:}

It is elementary to prove $E(x)\ll x^{1/2}$, and the prime number theorem implies $o(x^{1/2})$. The best known unconditional upper bound is of the shape $x^{1/2-o(1)}$, due to Walfisz \cite{Wa63}. Evelyn and Linfoot \cite{EvLi31} proved that\[E(x) \gg x^{1/4},\]and this is likely the true order of magnitude. The Riemann Hypothesis would follow from $E(x)\ll x^{1/4}$. 

The true order of magnitude is unknown even assuming the Riemann Hypothesis. Conditional on this assumption, the best known upper bound is\[E(x)\ll x^{\frac{11}{35}+o(1)},\]due to Liu \cite{Li16}.

\bigskip
\noindent\textbf{References:}

[EvLi31] Evelyn, C. J. A. and Linfoot, E. H., On a problem in the additive theory of numbers. Ann. of Math. (2) (1931), 261--270.
 
 [Li16] Liu, H.-Q., On the distribution of squarefree numbers. J. Number Theory (2016), 202--222.
 
 [Wa63] Walfisz, Arnold, Weylsche {E}xponentialsummen in der neueren {Z}ahlentheorie. (1963), 231.

\bigskip
\noindent\small{Source: \url{https://www.erdosproblems.com/969}}
\end{document}
